% 1_resumen.tex

\begin{abstract}
Este documento corresponde al reporte final del proyecto del curso IE0527: Ingeniería de Comunicaciones, el cual consiste en el desarrollo de un sistema de transmisión inalámbrica de audio.
Se integra el análisis del problema y la investigación teórica que fundamentan el diseño, la propuesta de arquitectura de hardware y el protocolo de transmisión, junto con la implementación de la solución y la descripción de las pruebas realizadas.
El sistema se compone de dos nodos con rol fijo, un transmisor y un receptor, basados en una Raspberry Pi Zero y un \textit{transceiver} NRF24L01+PA+LNA que opera en la banda ISM de 2.4 GHz, sin emplear protocolos inalámbricos de alto nivel.

Como fundamento, se investigaron los principios de digitalización de audio (frecuencia de muestreo, cuantización, formatos y esquemas de codificación) y las limitaciones del canal inalámbrico, considerando atenuación, interferencia, \textit{shadowing} y la capacidad máxima del canal según el teorema de Shannon.
A partir de ello, se seleccionó una configuración base de audio mono a 8 kHz con resolución de 8 bits por muestra en formato PCM, con ADPCM como optimización, y se definió un protocolo de aplicación de tres mensajes (\texttt{START}, \texttt{DATA} y \texttt{END}) sobre el mecanismo \textit{Enhanced ShockBurst} del NRF24L01, que gestiona la verificación CRC, las confirmaciones de recepción y las retransmisiones.
Finalmente, se documentan las verificaciones funcionales realizadas sobre la implementación en Python y el procedimiento de pruebas del sistema completo.
\end{abstract}

\begin{IEEEkeywords}
% En orden alfabético
ADPCM, audio digital, banda ISM, codificación de audio, digitalización, \textit{Enhanced ShockBurst}, fragmentación de paquetes, NRF24L01, Raspberry Pi Zero, transmisión inalámbrica.
\end{IEEEkeywords}
